\subsection{Unlearning Effectiveness}
\label{sec:effectiveness}
Tables~\ref{tab:iid-utility} and~\ref{tab:noniid-utility} compare the methods
under IID and Non-IID partitions. Values are mean $\pm$ sample SD across five
independent seeds (42--46). The target is class 0; the main method uses $rt=1$.
\begin{table}[H]\centering\small
\caption{IID utility and method time. Original time is pretraining cost; Retrain is the retained-data reference.}
\label{tab:iid-utility}
\begin{tabular}{lrrrrr}\toprule Method & $n$ & FA (\%) & RA (\%) & Acc. (\%) & Time (s)\\\midrule
Original & 5 & $79.08 \pm 0.37$ & $73.34 \pm 0.44$ & $73.92 \pm 0.40$ & $533.88 \pm 90.89$\\
Retrain & 5 & $0.00 \pm 0.00$ & $74.55 \pm 0.54$ & $67.10 \pm 0.48$ & $482.88 \pm 96.54$\\
Finetune & 5 & $78.12 \pm 0.40$ & $74.00 \pm 0.26$ & $74.41 \pm 0.23$ & $23.17 \pm 5.47$\\
GA & 5 & $17.38 \pm 1.75$ & $67.85 \pm 0.56$ & $62.80 \pm 0.51$ & $3.25 \pm 0.78$\\
Fed-EMNU & 5 & $6.80 \pm 1.35$ & $75.24 \pm 0.08$ & $68.40 \pm 0.20$ & $22.43 \pm 4.57$\\
UNSIR-style & 5 & $4.04 \pm 1.25$ & $74.80 \pm 0.22$ & $67.72 \pm 0.26$ & $13.34 \pm 2.56$\\
\bottomrule\end{tabular}\end{table}

\begin{table}[H]\centering\small
\caption{Non-IID ($\alpha=0.5$) utility and method time. Original time is pretraining cost; Retrain is the retained-data reference.}
\label{tab:noniid-utility}
\begin{tabular}{lrrrrr}\toprule Method & $n$ & FA (\%) & RA (\%) & Acc. (\%) & Time (s)\\\midrule
Original & 5 & $66.46 \pm 8.86$ & $69.62 \pm 1.10$ & $69.30 \pm 1.04$ & $631.98 \pm 371.26$\\
Retrain & 5 & $0.00 \pm 0.00$ & $70.05 \pm 1.01$ & $63.05 \pm 0.91$ & $802.48 \pm 842.82$\\
Finetune & 5 & $64.54 \pm 10.70$ & $71.59 \pm 0.86$ & $70.88 \pm 0.82$ & $48.12 \pm 59.14$\\
GA & 5 & $14.62 \pm 3.83$ & $66.38 \pm 2.86$ & $61.20 \pm 2.31$ & $4.82 \pm 3.30$\\
Fed-EMNU & 5 & $12.70 \pm 9.68$ & $72.16 \pm 0.70$ & $66.22 \pm 1.11$ & $40.41 \pm 41.25$\\
UNSIR-style & 5 & $4.36 \pm 3.39$ & $70.73 \pm 0.78$ & $64.09 \pm 0.79$ & $22.33 \pm 22.96$\\
\bottomrule\end{tabular}\end{table}
Fed-EMNU reduces target-class accuracy while retaining useful performance on
the other classes. Its mean RA exceeds that of Retrain and the UNSIR-style
adaptation in both partitions, but its FA remains above Retrain's zero and
above the UNSIR-style comparator. Thus the retained-utility result does not
establish equivalence to training without the target records.

The target attack AUC is $0.620\pm0.017$ (IID) and $0.596\pm0.006$ (Non-IID),
versus $0.508\pm0.026$ and $0.487\pm0.033$ for Retrain. Mean posterior TV is
$0.254\pm0.003$ and $0.313\pm0.008$, respectively. These diagnostics expose
remaining behavioral and membership differences. Complete reference distances,
target/retained attacks and controls are in Appendix~\ref{sec:core-details}.
The UNSIR-style comparator changes the noise penalty, noise steps and repair
budget together, preventing attribution of its differences to one component.

\textbf{Practical interpretation.} In the primary comparison, Fed-EMNU's mean
RA advantage over the UNSIR-style adaptation is 0.44 percentage points under
IID and 1.43 points under Non-IID. It costs more compute and retains more target
accuracy. Thus, the evidence does not establish a uniformly preferable method.
A deployment considering Fed-EMNU would need to value the retained-utility
gain, tolerate the measured residual target and membership behavior, and meet
its resource budget. For a requirement to verifiably remove individual training
influence, these experiments provide insufficient evidence to qualify the
method. Even the retraining reference is an empirical comparator rather than
a privacy certificate for the full system.

Figure~\ref{fig:additional-effectiveness} visualizes the main utility trade-off.

\begin{figure}[htbp]\centering
\includegraphics[width=\linewidth]{additional_effectiveness.pdf}
\caption{Primary-method target-class accuracy (FA) and retained accuracy (RA). Bars and whiskers are mean and sample SD across five seeds, using the same summaries as Tables~\ref{tab:iid-utility} and~\ref{tab:noniid-utility}. Lower FA indicates class suppression, not a deletion guarantee.}
\label{fig:additional-effectiveness}\end{figure}

\subsection{Unlearning Efficiency}
\label{sec:efficiency}
The paired Retrain/Fed-EMNU method-time ratio is $21.54\pm0.35$ under IID
and $19.56\pm0.71$ under Non-IID. The corresponding UNSIR-style/Fed-EMNU
ratios are $0.60\pm0.08$ and $0.55\pm0.03$, and GA/Fed-EMNU is
$0.14\pm0.01$ and $0.14\pm0.03$. A ratio below one means the named baseline
is faster. Finetune/Fed-EMNU ratios are $1.03\pm0.04$ and $1.06\pm0.16$.
These are paired compute/serialization ratios; evaluation and real transport
are excluded. Non-IID seed 45 took 114.20 s for Fed-EMNU, whereas the other
events took approximately 21--23 s. Its cause is unresolved and it is retained
in the reported mean and SD. The event carries 868.83 MiB of tensor payload
on average in both partitions; faster computation is not a network guarantee.

Figure~\ref{fig:additional-efficiency} shows the measured costs and variability.

\begin{figure}[htbp]\centering
\includegraphics[width=\linewidth]{additional_efficiency.pdf}
\caption{Primary method time, arithmetic mean across five seeds. Sample SD is reported in the main utility tables; it is omitted here because the Non-IID retraining SD exceeds its mean and a logarithmic error interval would cross zero. A logarithmic axis displays fast update methods and retained-data retraining together. Timing includes local computation and serialization, and excludes evaluation and real transport. The unresolved Non-IID timing outlier is retained.}
\label{fig:additional-efficiency}\end{figure}

\subsection{Ablation of Unlearning and Enhancement}
\label{sec:ablation}
To explain how the existing components contribute, Figure~\ref{fig:ablation}
compares the full workflow with unlearning alone (impair-only) and enhancement
alone (repair-only), using the same request checkpoints across five seeds.
These controls remove one existing stage and retain the prescribed budget of
the remaining stage; they introduce no additional component.

\textbf{Role of enhancement.} With unlearning alone, IID FA is
$1.92\pm1.29$\% but RA falls to $72.61\pm1.51$\%, compared with
$6.80\pm1.35$\% FA and $75.24\pm0.08$\% RA for the full workflow.
Under Non-IID, RA rises from $36.78\pm22.04$\% without enhancement to
$72.16\pm0.70$\% with both stages, while FA rises from $1.08\pm1.32$\%
to $12.70\pm9.68$\%. Enhancement therefore recovers retained utility in
this recipe, with some recovery of target-class predictions as well.

\textbf{Role of target-directed unlearning.} Enhancement alone preserves
RA ($75.56\pm0.33$\% IID; $72.82\pm0.76$\% Non-IID) but leaves FA at
$41.42\pm1.83$\% and $41.00\pm7.45$\%, respectively. Retained-only
enhancement is therefore insufficient to obtain the full workflow's target
suppression. Together, these controls support the complementary functions of
the two stages: target-directed updating supplies suppression, and enhancement
recovers retained-task utility. They explain the design's measured trade-off,
rather than establish that either primitive is new. Removing a stage also
reduces computation, so these are stage-removal controls, not equal-compute
comparisons or a proof of record deletion. They do not separately isolate
optimized noise from arbitrary synthetic inputs or other impair objectives.
\begin{figure}[htbp]\centering
\includegraphics[width=\linewidth]{fig5_ablation_revised.pdf}
\caption{Two-stage ablation under IID and Non-IID partitions.
Bars show mean Acc., FA and RA; whiskers show sample SD across five seeds.
Forget-only denotes impair-only; Enhance-only denotes repair-only.
All controls use the same request checkpoints and retained-data references.}
\label{fig:ablation}\end{figure}

\subsection{Local Retained-Data Selection}
\label{sec:selection-results}
Figure~\ref{fig:filtering} compares random and confidence-based selection. Both controls use $rt=0.5$, identical per-class counts, and paired
checkpoints. Confidence minus random selection changes RA by
$-1.50\pm0.27$ percentage points under IID and $-0.84\pm0.61$ under Non-IID.
Their descriptive 95\% intervals are $[-1.837,-1.154]$ and
$[-1.592,-0.083]$ (Appendix~\ref{sec:uncertainty}). The evaluated recipe does
not support a confidence-selection advantage. These results do not establish an optimal retained ratio: the selection
comparison uses a single ratio, and the $rt=1$ main recipe is not a matched
ratio-only control for these experiments.
\begin{figure}[htbp]\centering
\includegraphics[width=\linewidth]{fig6_filtering_revised.pdf}
\caption{Matched retained-data selection at $rt=0.5$. Both partitions are evaluated. Bars and whiskers show mean and sample
SD across five paired seeds; metric colors match Figure~\ref{fig:ablation}.}
\label{fig:filtering}\end{figure}

\subsection{Random-Seed Variation}
\label{sec:seed-results}
Figure~\ref{fig:seed} shows mean and sample SD across five seeds in both partitions. Non-IID FA varies substantially ($12.70\pm9.68$\%),
although RA varies less ($72.16\pm0.70$\%). Variation is therefore not uniformly small across metrics and partitions. Primary confidence intervals and
reference-to-reference variation are reported in Appendices~\ref{sec:uncertainty}
and~\ref{sec:core-details}; small samples and variation in partition and schedule
limit general stability claims.
\begin{figure}[htbp]\centering
\includegraphics[width=\linewidth]{fig8_seed_revised.pdf}
\caption{Seed variation of the full method: mean $\pm$ sample SD across seeds
42--46. Both partitions are shown. Low FA alone is not a privacy or deletion
guarantee; the attack and reference diagnostics must be considered jointly.}
\label{fig:seed}\end{figure}

\subsection{Additional Evaluation}
\label{sec:extension-summary}
The complete additional studies are retained in Appendix~\ref{sec:extensions}.
ResNet34 extends the architecture evidence within the ResNet family. Stronger
heterogeneity ($\alpha=0.1$) exposes poor starting utility and sparse target
ownership, including the aggregation limitation in Section~\ref{sec:protocol}.
In CIFAR-100, FA remains $33.00\pm12.49$\% (IID) and $24.67\pm13.32$\%
(Non-IID), and target AUC exceeds the retraining reference; the fixed recipe
does not establish effective cross-dataset deletion.

The six PAR-code-based events use the same three-seed starting checkpoints as
their core comparators, with the adaptation described in
Appendix~\ref{sec:parprotocol}. PAR adaptation FA is $63.00\pm2.70$\% (IID)
and $54.23\pm5.81$\%; RA is $74.76\pm0.33$\% and $72.14\pm0.47$\%.
Its native budgets differ, and unavailable full-paper verification prevents
claiming an exact reproduction or refutation of the published method.

Twenty-client tests yield RA $71.41\pm0.80$\% (IID) and $66.69\pm1.41$\%
(Non-IID), with paired 20-minus-10-client RA differences of
$-3.88\pm0.84$ and $-5.46\pm0.58$ percentage points. Partitioning and local
batching also change, so this is a limited scale test. The two-request study
observes no class-0 rebound in six Fed-EMNU paths, with variable class-1
suppression. Five retained-only continuation rounds after the single-class
event also show no FA rebound in six paths, but target membership scores
remain different from the continued retained-data reference. Neither study
validates arbitrary request streams or future training. Full tables and
protocol boundaries remain available in the appendices.


\subsection{Controlled Comparisons and Robust Timing}
\label{sec:controlled-summary}
Matched objective controls use three paired seeds, common checkpoints, full
retained data and identical random ordering. At 80 noise steps/three repair
epochs, Fed-minus-UNSIR RA is $0.063\pm0.371$ percentage points under IID
and $-0.463\pm0.361$ under Non-IID; at 40 steps/one repair epoch it is
$-0.004\pm0.306$ and $-0.356\pm0.222$ points. These controls do not establish
a consistent advantage of the Fed-EMNU penalty. Optimized noise lowers target
attack AUC relative to Gaussian inputs by $0.029\pm0.012$ and
$0.022\pm0.018$, respectively, but has mixed accuracy effects.
Retained-count event weighting raises Non-IID RA by $0.896\pm0.224$ points
and FA by $0.600\pm0.557$ relative to equal weighting.

Separate comparisons impose common maximum compute/serialization caps of
10, 20, 30 and 40 seconds. All 48 events finish within their caps. Fed-EMNU
has higher mean RA at every cap, by 1.41/0.49/0.48/0.39 points under IID
and 4.50/1.81/1.29/1.56 points under Non-IID, but also higher FA and target
attack AUC throughout. These are time-capped adaptations with reported
realized steps and consumed times; early completion is not padded. The results
show a utility--suppression trade-off rather than dominance or a privacy gain.

Robust summaries retain the historical Non-IID timing outlier. Six technical
repeats reproduce the saved predictions exactly but do not reproduce its large
delay: seed 45 takes a median 25.76 seconds, versus 114.20 seconds originally.
The historical cause remains unresolved; repeats are not additional scientific
seeds. Appendix~\ref{sec:controlled-details} gives the full control protocols,
tables, timing distributions and limits, separately from the five-seed primary
comparison.

\subsection{Applicability and Failure Conditions}
\label{sec:applicability}
The fixed recipe should not be extrapolated beyond the evaluated conditions.
On CIFAR-100, target accuracy remains 33.00\% under IID and 24.67\% under
Non-IID, with target membership AUC above the matching references. The method
therefore fails to establish reliable class revocation on that dataset under
the stated recipe. At twenty clients, retained accuracy drops by 3.88 and 5.46
percentage points relative to paired ten-client runs. This does not support
a claim that utility is maintained as the federation grows. With stronger
heterogeneity, poor starting utility and sparse target ownership expose the
limits of updating only target holders; the current aggregation rule gives
other clients no separate repair contribution. ResNet34 adds depth within one
architecture family, without establishing architecture independence.

The mechanism is consequently an experimental option for approximate
class suppression in cooperative, available-client settings with similar
data and training conditions. Deployment would require representative local
validation of target behavior, retained utility, membership exposure and
resource use. The observed failures are part of the method's scope, not
exceptions removed from the evaluation. No universal operating threshold or
certified-deletion use case is established by this study.
